\section{The fractional-shift bridge}\label{sec:bridge}\label{sec:bridge-precision}

This section builds the fractional-shift bridge, the steps \textsf{Spread}, \textsf{TensorFFT}, and \textsf{Deconvolve}
of Algorithm~\ref{alg:main} (Section~\ref{sec:overview-five-steps}), which reduces a one-dimensional discrete Fourier transform $F_R$
to a $d$-dimensional tensor transform $F_{\mathbf t}$ on the same power-of-two grid $R=t_1\cdots t_d$. Before the tensor
transform, each axis above the first is smoothed by a shifted complex sinc-Gaussian filter; afterwards, for each
smoothed axis, a two-term cyclic recurrence in the frequency domain undoes that smoothing to working precision. For
Theorem~\ref{thm:main}, this reduction avoids prime transform lengths and grid expansion, and sets the linear
working-precision budget $P=O(p+\lg R)$ used throughout the multiplier. A first-time reader willing to treat the bridge
as a black box needs only Definition~\ref{def:middle-contract} for the middle-transform contract,
Theorem~\ref{thm:bridge-precision} for the bridge precision contract, and Lemma~\ref{lem:bridge-D} for the axis-exposure
bound, and may skip directly to Section~\ref{sec:walsh}.

Sections~\ref{sec:bridge-conventions}--\ref{sec:bridge-B} carry out the continuous and exact Fourier analysis. We spread
the input along axes $a=d,d-1,\dots,2$ in descending order using a negative half-bin phase carrier, and the triangular
alias identity (Lemma~\ref{lem:bridge-A}) expresses the tensor DFT of the spread array as an ordered product
$\cA_2\cdots\cA_d F_R x$ of non-commuting alias operators. Lemma~\ref{lem:bridge-B} then approximates each $\cA_a$ to
within $0.120\cdot 2^{-\tw}$ by a two-term cyclic operator $B_a$ whose inverse has $\ell^\infty$ norm at most
$4t_a/7\le t_a$, uniformly in its cycle length $A_a$. Applying $B_2^{-1},\dots,B_d^{-1}$ in that order therefore
recovers $F_R x$ up to the omitted-alias error bounded in~\eqref{eq:lemma-B-total}. Two converse results in
Appendix~\ref{app:deferred}, Propositions~\ref{prop:positive-fails} and~\ref{prop:real-kernel-fails}, show that within
this two-alias framework both the negative sign of the half-bin carrier and a complex kernel are required for uniform
bounds.

Sections~\ref{sec:bridge-alg}--\ref{sec:bridge-C} turn this analytic reduction into a fixed-point algorithm and prove
its error and bit-width bounds in Lemma~\ref{lem:bridge-C} and Theorem~\ref{thm:bridge-precision}. The rational-series
evaluation of the filter and phase tables is deferred to Lemma~\ref{lem:phase-tables} in Appendix~\ref{app:deferred}.
Finally, Section~\ref{sec:bridge-D} gives the $O(d)$ direct axis-exposure schedule on multitape machines in
Lemma~\ref{lem:bridge-D}.

\subsection{Index conventions and the carrier}\label{sec:bridge-conventions}

This subsection fixes the mixed-radix indexing of inputs and frequencies on the tensor grid and defines the shifted
complex filter used on each spread axis.

\paragraph{Parameters and index conventions.}
Let $R=t_1\cdots t_d=2^k$ with every $t_i\ge 2$ a power of two, where $k$ denotes $\lg R$ throughout the paper. Put
$A_a=\prod_{i<a}t_i$ for $1\le a\le d+1$, so that $A_1=1$ and $A_{d+1}=R$. We write input indices as
$j=\sum_{i=1}^d j_i A_i$ with $0\le j_i<t_i$, making $j_1$ the least significant digit, and identify $x\in\C^R$ with
the tensor $x_{j_1,\dots,j_d}$. Frequencies use the reversed digit order:
\begin{equation}\label{eq:freq-digits}
  f(\mathbf u)=\sum_{i=1}^d u_i\,\frac{R}{A_{i+1}},\qquad 0\le u_i<t_i,
\end{equation}
so $u_d$ is the least significant digit. The map $\mathbf u\mapsto f(\mathbf u)$ is a bijection onto $[R]$, and we write
$u_a(f)$ for the $a$-th frequency digit of $f\in[R]$. For $2\le a\le d$, the \emph{fractional shift} of axis~$a$ is
\begin{equation}\label{eq:delta}
  \delta_a(\mathbf j)=\frac{\sum_{i<a}j_i A_i}{A_a}\in[0,1),
\end{equation}
which depends only on the input digits below~$a$. In a multi-axis Cooley--Tukey decomposition of $F_R$, the twiddle
factor coupling axis~$a$ to the lower axes acts as a fractional shift by $\delta_a(\mathbf j)$ along axis~$a$,
generalizing the two-dimensional observation of~\cite{MS}.

As in~\cite{MS}, our base smoothing kernel is a sinc-Gaussian whose Fourier transform $G$ approximates the indicator of
$[-\tfrac12,\tfrac12]$. On each spread axis $a\ge 2$, we multiply this base kernel by a negative half-bin phase so that
frequency samples move from the grid points $u/t_a$ to the midpoints $(u+\tfrac12)/t_a$. The next definition records
this modulated carrier and its $t_a$-periodic folded filter, which supply the weights for axis spreading.

\begin{definition}[Kernel and carrier]\label{def:carrier}
Let $\tw\ge 8$ be the alias width, and fix $\alpha$ satisfying
\begin{equation}\label{eq:alpha}
  2\sqrt{\tw\ln 2/\pi}\le\alpha\le 4\sqrt{\tw\ln 2/\pi},
\end{equation}
with default $\alpha=2\sqrt{\tw\ln 2/\pi}$. Put $\alpha_\pi=\sqrt\pi\,\alpha$, so $\alpha_\pi\ge 2\sqrt{\tw\ln 2}$. With
$\sinc(s)=\sin(\pi s)/(\pi s)$ for $s\ne 0$, $\sinc(0)=1$, and Fourier transform $\hat\phi(\xi)=\int_{\R}\phi(s)e^{-2\pi i\xi s}\,ds$,
the transform of $\sinc(s)\,e^{-\pi s^2/\alpha^2}$ is
\begin{equation}\label{eq:G}
  G(\xi)=\int_{\xi-1/2}^{\xi+1/2}\alpha e^{-\pi\alpha^2 s^2}\,ds
  =\tfrac12\bigl[\erf\bigl(\alpha_\pi(\xi+\tfrac12)\bigr)-\erf\bigl(\alpha_\pi(\xi-\tfrac12)\bigr)\bigr].
\end{equation}
On axis $a\in\{2,\dots,d\}$, the \emph{negative half-bin carrier} is
\begin{equation}\label{eq:carrier}
  g_a(s)=e^{-\pi i s/t_a}\,\sinc(s)\,e^{-\pi s^2/\alpha^2},\qquad
  \hat g_a(\xi)=G\Bigl(\xi+\frac{1}{2t_a}\Bigr),
\end{equation}
and its \emph{folded filter} is the $t_a$-periodic function of $e\in\Z$ given by
\begin{equation}\label{eq:folded}
  H^{(a)}_\delta(e)=\sum_{s\in e+t_a\Z}g_a(s+\delta).
\end{equation}
\end{definition}

Because the carrier phase $e^{-\pi i s/t_a}$ is antiperiodic of period $t_a$ rather than periodic, the modulated kernel
itself must be folded as in~\eqref{eq:folded}. Since $e^{-\pi i(s+\delta)/t_a}=(-1)^{(e-s)/t_a}e^{-\pi i(e+\delta)/t_a}$
for $s\in e+t_a\Z$,
\begin{equation}\label{eq:negacyclic}
  H^{(a)}_\delta(e)=e^{-\pi i(e+\delta)/t_a}\sum_{s\in e+t_a\Z}(-1)^{(e-s)/t_a}\,g_0(s+\delta),\qquad
  g_0(s)=\sinc(s)\,e^{-\pi s^2/\alpha^2}.
\end{equation}
Periodizing $g_0$ first and multiplying by the carrier afterwards would drop the alternating signs $(-1)^{(e-s)/t_a}$.
Identity~\eqref{eq:negacyclic} also shows that $H^{(a)}_\delta$ can be applied as a real negacyclic convolution between
two diagonal phase scans. In Section~\ref{sec:bridge-alg}, our algorithm evaluates the complex filter directly.

\subsection{Triangular alias identity}\label{sec:bridge-A}

This subsection shows that spreading the axes in descending order expresses the tensor transform of the spread array as
an exact ordered product of one-axis alias operators applied to $F_R x$.

We spread one axis at a time. For $2\le a\le d$, let $\mathcal S_a$ spread axis~$a$ while holding all other coordinates
fixed:
\begin{equation}\label{eq:spread-axis}
  (\mathcal S_a V)_{\dots,k_a,\dots}=\sum_{j_a\in[t_a]}V_{\dots,j_a,\dots}\,H^{(a)}_{\delta_a(\mathbf j)}(j_a-k_a),
\end{equation}
where $\delta_a(\mathbf j)$ is read from the coordinates below~$a$ of the current array. We apply these operators in the
descending order $a=d,d-1,\dots,2$ and set $U=\mathcal S_2\mathcal S_3\cdots\mathcal S_d x$. When axis~$a$ is spread,
coordinates $1,\dots,a-1$ are still the original input digits $j_1,\dots,j_{a-1}$, because they are spread only at later
stages. Hence
\begin{equation}\label{eq:spread-product}
  U_{j_1,k_2,\dots,k_d}=\sum_{j_2,\dots,j_d}x_{\mathbf j}\prod_{a=2}^d H^{(a)}_{\delta_a(\mathbf j)}(j_a-k_a).
\end{equation}
In the opposite order, $\delta_a$ would be evaluated after lower digits had already been replaced by spread indices, and
\eqref{eq:spread-product} would fail.

For $2\le a\le d$, the \emph{alias operator} $\cA_a$ on $\C^R$ collects the frequency aliases $f-\nu R/A_a$ for integer
shifts $\nu\in\Z$ generated by Poisson summation along axis~$a$, weighted by half-bin samples of~$G$:
\begin{equation}\label{eq:alias-op}
  (\cA_a Z)_f=\sum_{\nu\in\Z}G\Bigl(\frac{u_a(f)+1/2}{t_a}-\nu\Bigr)\,Z_{(f-\nu R/A_a)\bmod R}.
\end{equation}
The weights of $\cA_a$ depend only on the digit $u_a(f)$, whereas shifting $f$ by a multiple of $R/A_a$ alters only the
more significant digits $u_1,\dots,u_{a-1}$.

In the product $\cA_2\cA_3\cdots\cA_d$, each factor shifts only frequency digits that no factor to its right reads. The
next lemma uses this triangular dependency to express the tensor transform of the spread array $U$ directly in terms of
the target transform $Z=F_R x$. Because the alias operators do not commute, their order dictates the left-to-right order
in which deconvolution later inverts their truncations.

\begin{lemma}[Triangular alias identity]\label{lem:bridge-A}
With $U=\mathcal S_2\mathcal S_3\cdots\mathcal S_d x$ and $Z=F_R x$,
\begin{equation}\label{eq:lemma-A}
  F_{\mathbf t}U=\cA_2\cA_3\cdots\cA_d\,Z,
\end{equation}
where the operators compose from right to left. Explicitly, for every $\mathbf u\in[t_1]\times\cdots\times[t_d]$,
\begin{equation}\label{eq:lemma-A-explicit}
  (F_{\mathbf t}U)_{\mathbf u}=\sum_{\boldsymbol\nu\in\Z^{d-1}}
  \Bigl[\prod_{a=2}^d G\Bigl(\frac{u_a+1/2}{t_a}-\nu_a\Bigr)\Bigr]\,
  Z_{(f(\mathbf u)-\sum_{a=2}^d\nu_a R/A_a)\bmod R},
\end{equation}
where $\boldsymbol\nu=(\nu_2,\dots,\nu_d)$. The operators $\cA_a$ do not commute in general, and~\eqref{eq:lemma-A}
holds in the stated order.
\end{lemma}

\emph{Proof idea.} We unfold the periodic filter on each spread axis and apply Poisson summation to express
$(F_{\mathbf t}U)_{\mathbf u}$ as an absolutely convergent sum over alias vectors $\boldsymbol\nu\in\Z^{d-1}$. Comparing
mixed-radix phases modulo~$1$ identifies each inner sum over the input grid with a shifted entry of $Z=F_R x$, yielding
\eqref{eq:lemma-A-explicit}. Expanding $\cA_2\cA_3\cdots\cA_d Z$ from left to right then shows that earlier frequency
shifts never modify the frequency digit read by a later operator.

\begin{proof}
\emph{Explicit formula.} For any Schwartz function $\phi$ and real $z,\xi$, applying Poisson summation to
$s\mapsto\phi(z-s)e^{2\pi i\xi s}$ yields
\begin{equation}\label{eq:poisson}
  \sum_{e\in\Z}\phi(z-e)e^{2\pi i e\xi}=\sum_{\nu\in\Z}\hat\phi(\xi-\nu)\,e^{2\pi i z(\xi-\nu)}.
\end{equation}
No symmetry of $\phi$ is required, so~\eqref{eq:poisson} applies directly to the complex carriers $g_a$. Substitute
\eqref{eq:spread-product} into $F_{\mathbf t}U$. For fixed $\mathbf j$ and $a\in\{2,\dots,d\}$, unfolding the
$t_a$-periodic filter $H^{(a)}_{\delta_a}$ and applying~\eqref{eq:poisson} with $\phi=g_a$, $z=j_a+\delta_a(\mathbf j)$,
and $\xi=u_a/t_a$ gives
\[
  \sum_{k_a\in[t_a]}H^{(a)}_{\delta_a}(j_a-k_a)\,e^{2\pi i k_a u_a/t_a}
  =\sum_{\nu_a\in\Z}\hat g_a\Bigl(\frac{u_a}{t_a}-\nu_a\Bigr)\,e^{2\pi i(j_a+\delta_a)(u_a/t_a-\nu_a)}.
\]
By the Gaussian decay of $g_a$ and $\hat g_a$, all sums converge absolutely and may be interchanged:
\[
  (F_{\mathbf t}U)_{\mathbf u}=\sum_{\boldsymbol\nu\in\Z^{d-1}}
  \Bigl[\prod_{a=2}^d\hat g_a\Bigl(\frac{u_a}{t_a}-\nu_a\Bigr)\Bigr]\,
  \frac1R\sum_{\mathbf j}x_{\mathbf j}
  \exp\Bigl(2\pi i\Bigl[\frac{j_1 u_1}{t_1}+\sum_{a=2}^d(j_a+\delta_a)\Bigl(\frac{u_a}{t_a}-\nu_a\Bigr)\Bigr]\Bigr).
\]
Since $A_i/A_{a+1}\in\Z$ for $i>a$, we have $j/A_{a+1}\equiv(j_a+\delta_a)/t_a\pmod 1$ and $j/A_a\equiv\delta_a\pmod 1$.
Moreover, $j/t_1\equiv j_1/t_1\pmod 1$ and $\nu_a j_a\in\Z$. The phase in brackets is therefore congruent modulo~$1$ to
\[
  \frac{j}{R}\Bigl(\sum_{i=1}^d u_i\,\frac{R}{A_{i+1}}-\sum_{a=2}^d\nu_a\,\frac{R}{A_a}\Bigr),
\]
so the inner normalized sum over $\mathbf j$ equals $Z_{(f(\mathbf u)-\sum_{a=2}^d\nu_a R/A_a)\bmod R}$. Substituting
$\hat g_a(u_a/t_a-\nu_a)=G((u_a+\tfrac12)/t_a-\nu_a)$ yields~\eqref{eq:lemma-A-explicit}.

\emph{Product form.} Subtracting $\nu R/A_a$ from $f$ modifies only the digits $u_1,\dots,u_{a-1}$, because $R/A_a$ is
the positional weight of $u_{a-1}$ in~\eqref{eq:freq-digits} and borrows propagate toward $u_1$. The lower digits
$u_a,\dots,u_d$ remain unchanged. When $(\cA_2\cA_3\cdots\cA_d Z)_{f(\mathbf u)}$ is expanded from left to right, every
frequency shift accumulated before the weight of $\cA_a$ is evaluated comes from factors $\cA_{a'}$ with $a'<a$, which
leave the digit $u_a$ invariant. Every weight in the expansion is therefore evaluated at the original digit $u_a$ of
$f(\mathbf u)$, matching the right-hand side of~\eqref{eq:lemma-A-explicit}. In the reverse order a shift from
$\cA_{a+1}$ can alter $u_a$ before $\cA_a$ reads it, so the factors do not commute.
\end{proof}

Consequently, deconvolution must invert the approximate factors from left to right: first $B_2^{-1}$, then $B_3^{-1}$,
and so on through $B_d^{-1}$, where $B_a$ is the two-alias truncation of $\cA_a$ analyzed in
Section~\ref{sec:bridge-B}. The next remark illustrates why reversing this order fails already in dimension three, the
smallest case with two alias operators.

\begin{remark}[Non-commutativity in dimension three]\label{rem:bridge-d3}
Take $\mathbf t=(2,2,2)$, so $R=8$, $A_2=2$, $A_3=4$, $j=j_1+2j_2+4j_3$, and $f=4u_1+2u_2+u_3$, with retained aliases
$f-4\nu_2-2\nu_3$ ($\nu_2,\nu_3\in\{0,1\}$). For $x=e_3$ ($(j_1,j_2,j_3)=(1,1,0)$), we have $\delta_2=\tfrac12$,
$\delta_3=\tfrac34$, and $Z_f=\tfrac18 e^{2\pi i\,3f/8}$. At $\mathbf u=(1,0,1)$ ($f=5$), the four retained aliases
$(\nu_2,\nu_3)=(0,0),(1,0),(0,1),(1,1)$ land at frequencies $5,1,3,7$, with values $1,-1,i,-i$ times $Z_5$; with
$c_0=G(\tfrac14)$ and $c_1=G(\tfrac34)$, axis~$2$ contributes weights $(c_0,c_1)$ and axis~$3$ weights $(c_1,c_0)$, so
up to omitted aliases~\eqref{eq:lemma-A-explicit} gives
$\frac{c_0-c_1}{8\sqrt2}\bigl((c_0+c_1)+i(c_0-c_1)\bigr)$.
Similarly, writing $c_{a,\nu}(u)=G((u+\tfrac12)/t_a-\nu)$ for the weights of the two-term operator $B_a$
in~\eqref{eq:Bh} below and taking $Z$ to be the unit vector at frequency~$6$, evaluating $B_2 B_3 Z$ and $B_3 B_2 Z$ at
frequency~$0$ gives $c_{2,0}(0)c_{3,1}(0)$ and $c_{3,1}(0)c_{2,0}(1)$, which differ since
$c_{2,0}(0)=G(\tfrac14)>G(\tfrac34)=c_{2,0}(1)$ and $c_{3,1}(0)>0$.
\end{remark}

\subsection{Two-term cyclic deconvolution}\label{sec:bridge-B}

This subsection replaces each infinite alias operator $\cA_a$ by a two-term cyclic operator $B_a$ and proves that $B_a$
is well conditioned uniformly in its cycle length.

Because each $\cA_a$ sums over all $\nu\in\Z$, the algorithm inverts its two-alias truncation. For $2\le a\le d$, the
\emph{two-term cyclic operator} $B_a$ retains the terms $\nu\in\{0,1\}$ of~\eqref{eq:alias-op}: with
$y=(u_a(f)+\tfrac12)/t_a\in(0,1)$,
\begin{equation}\label{eq:Bh}
  (B_a Z)_f=c_0 Z_f+c_1 Z_{(f-R/A_a)\bmod R},\qquad c_0=G(y),\quad c_1=G(y-1).
\end{equation}
Each orbit of $f\mapsto(f-R/A_a)\bmod R$ has length $A_a$, and the digit $u_a(f)$ is constant along every orbit. Thus
$B_a$ decomposes as a direct sum of cyclic bidiagonal operators $B\colon Z_j\mapsto c_0 Z_j+c_1 Z_{j-1}$ on cycles of
length $L=A_a$, each with constant coefficients $(c_0,c_1)$. In the balanced shape of Section~\ref{sec:bridge-alg},
$A_a$ grows up to $R/r$, far larger than $t_a$, so the conditioning of $B^{-1}$ must be independent of the cycle length
$L$.

Because the negative half-bin carrier places every sample $y=(u+\tfrac12)/t_a$ strictly inside $(0,1)$, all omitted
aliases lie at distance greater than $\tfrac12$ from the Gaussian center. Since $t_a$ is even, no sample hits the
Nyquist midpoint $y=\tfrac12$ where $c_0$ and $c_1$ would tie. The next lemma turns these two geometric margins into a
uniform $O(2^{-\tw})$ truncation bound on $\cA_a-B_a$ and a cycle-length-independent bound
$\|B_a^{-1}\|_\infty\le 4t_a/7$, ensuring that inverting $B_2,\dots,B_d$ recovers $F_R x$ up to the omitted-alias error
of order $(d-1)(R/t_1)2^{-\tw}\|x\|_\infty$ in~\eqref{eq:lemma-B-total}.

\begin{lemma}[Negative half-bin kernel and conditioning]\label{lem:bridge-B}
Let $\tw\ge 8$, let $\alpha$ satisfy~\eqref{eq:alpha}, and let $t=t_a\ge 2$ be a power of two ($2\le a\le d$). For every
frequency digit $0\le u<t$, with $y=(u+\tfrac12)/t$, $c_0=G(y)$, and $c_1=G(y-1)$:
\begin{enumerate}[label=(\roman*),leftmargin=*]
  \item $\|\cA_a\|_\infty\le 1$ and $\|B_a\|_\infty\le 1$;
  \item \textup{(omitted aliases)} $\|\cA_a-B_a\|_\infty<0.120\cdot 2^{-\tw}$;
  \item \textup{(passband)} $c_{\max}:=\max(c_0,c_1)\ge\tfrac12-0.060\cdot 2^{-4\tw}>0.49999999$;
  \item \textup{(separation)} $|c_0-c_1|\ge 7/(4t)$;
  \item \textup{(conditioning)} on a cycle of any length $L\ge 1$, the operator $B\colon Z_j\mapsto c_0 Z_j+c_1 Z_{j-1}$
    is invertible with $\|B^{-1}\|_\infty\le 1/|c_0-c_1|\le 4t/7$ (with equality in the first bound when $L$ is even);
    consequently $\|B_a^{-1}\|_\infty\le 4t_a/7\le t_a$;
  \item \textup{(product)} $\|(B_2\cdots B_d)^{-1}\|_\infty\le(4/7)^{d-1}R/t_1\le R/t_1$, and for $Y=F_{\mathbf t}U$,
    \begin{equation}\label{eq:lemma-B-total}
      \|(B_2\cdots B_d)^{-1}Y-F_R x\|_\infty\le 0.120\,(d-1)\,\frac{R}{t_1}\,2^{-\tw}\,\|x\|_\infty.
    \end{equation}
\end{enumerate}
\end{lemma}

\emph{Proof idea.} Parts~(i)--(iii) bound the row sums, omitted Gaussian tails, and passband weight of $G$ by tiling
$\R$ with unit intervals and using the convexity of the complementary error function. Part~(iv) lower-bounds the weight
gap $|c_0-c_1|$ by differentiating $G(\tfrac12-e)-G(\tfrac12+e)$ and using the half-bin distance $|y-\tfrac12|\ge 1/(2t)$
from the midpoint. Parts~(v)--(vi) invert each cyclic bidiagonal block via a finite Neumann series whose row sum cancels
the cycle-closure denominator, and then telescope the product error across axes.

\begin{proof}
\emph{(i) Norms.} The intervals $[\xi-\nu-\tfrac12,\xi-\nu+\tfrac12]$ for $\nu\in\Z$ tile $\R$, so
$\sum_{\nu\in\Z}G(\xi-\nu)=\int_{\R}\alpha e^{-\pi\alpha^2 s^2}\,ds=1$ with all terms positive. The absolute row sums of
$\cA_a$ and $B_a$ are at most $\sum_{\nu\in\Z}G(y-\nu)=1$ and $c_0+c_1<1$, respectively.

\emph{(ii) Omitted aliases.} Removing the intervals for $\nu\in\{0,1\}$ leaves
$(-\infty,y-\tfrac32]\cup[y+\tfrac12,\infty)$, so the omitted alias weight at $y$ is
$\Psi(y)=\tfrac12\erfc(\alpha_\pi(y+\tfrac12))+\tfrac12\erfc(\alpha_\pi(\tfrac32-y))$, which bounds every absolute row
sum of $\cA_a-B_a$. Since $\erfc''>0$ on $(0,\infty)$, $\Psi$ is convex on $[0,1]$ and symmetric about $\tfrac12$, whence
$\Psi(y)\le\Psi(0)$ for all $y\in(0,1)$. Using $\erfc z\le e^{-z^2}/(\sqrt\pi z)$ for $z>0$ and
$\alpha_\pi/2\ge\sqrt{\tw\ln 2}$, we obtain for $\tw\ge 8$ that
\[
  \Psi(0)\le\frac{2^{-\tw}}{2\sqrt{\pi\tw\ln 2}}+\frac{2^{-9\tw}}{6\sqrt{\pi\tw\ln 2}}<0.120\cdot 2^{-\tw}.
\]
Because every half-bin sample $y=(u+\tfrac12)/t$ lies in $(0,1)$, this bound is uniform in $t\ge 2$.

\emph{(iii) Passband.} Since $G$ is even and strictly decreasing on $[0,\infty)$ and $\min(y,1-y)\le\tfrac12$, we have
$c_{\max}\ge G(\tfrac12)=\tfrac12\erf(\alpha_\pi)=\tfrac12-\tfrac12\erfc(\alpha_\pi)$, and
$\tfrac12\erfc(\alpha_\pi)\le e^{-\alpha_\pi^2}/(2\sqrt\pi\,\alpha_\pi)<0.060\cdot 2^{-4\tw}$ for $\tw\ge 8$.

\emph{(iv) Separation.} For $0\le e\le\tfrac12$, put $\Delta_G(e)=G(\tfrac12-e)-G(\tfrac12+e)$. Differentiating
\eqref{eq:G} and using $\pi\alpha^2=\alpha_\pi^2$ gives
\[
  \Delta_G'(e)=\alpha\bigl[2e^{-\alpha_\pi^2 e^2}-e^{-\alpha_\pi^2(1-e)^2}-e^{-\alpha_\pi^2(1+e)^2}\bigr]
  =2\alpha e^{-\alpha_\pi^2 e^2}\bigl(1-e^{-\alpha_\pi^2}\cosh(2\alpha_\pi^2 e)\bigr)>0
\]
on $[0,\tfrac12]$, so $\Delta_G$ is strictly increasing on $[0,\tfrac12]$. For $0\le e\le\tfrac14$, we have
$e^{-\alpha_\pi^2}\cosh(2\alpha_\pi^2 e)\le e^{-\alpha_\pi^2(1-2e)}\le e^{-\alpha_\pi^2/2}\le 2^{-2\tw}\le 2^{-16}$, so
$\Delta_G'(e)\ge 2\alpha(1-2^{-16})e^{-\alpha_\pi^2 e^2}$, and integrating from $0$ yields
$\Delta_G(e)\ge(1-2^{-16})\erf(\alpha_\pi e)$. Because $t\ge 2$ is even, the midpoint $(t-1)/2$ is a half-integer, so
every $u\in[t]$ satisfies $e=|y-\tfrac12|=|u-\tfrac{t-1}{2}|/t\ge\frac{1}{2t}$, with $\frac{1}{2t}\le\tfrac14$. By
concavity of $\erf$ on $[0,\infty)$ and $2/t\le 1$,
\[
  \erf\Bigl(\frac{\alpha_\pi}{2t}\Bigr)\ge\frac{2}{t}\,\erf\Bigl(\frac{\alpha_\pi}{4}\Bigr)\ge\frac{2}{t}\,\erf(\sqrt{2\ln 2}),
\]
since $\alpha_\pi/4\ge\tfrac12\sqrt{\tw\ln 2}\ge\sqrt{2\ln 2}$. The tail bound of part~(iii) gives
$\erf(\sqrt{2\ln 2})>1-\frac{1}{4\sqrt{2\pi\ln 2}}>\frac{22}{25}$, and therefore
\[
  t\,|c_0-c_1|\ge t\,\Delta_G\Bigl(\frac{1}{2t}\Bigr)>2\,(1-2^{-16})\,\frac{22}{25}>\frac{7}{4}.
\]

\emph{(v) Conditioning.} Set $\rho=\min(c_0,c_1)/c_{\max}<1$ and let $T$ be the cyclic shift $(TZ)_j=Z_{(j-1)\bmod L}$.
If $c_0>c_1$, then $BZ=Y$ reads $(I+\rho T)Z=Y/c_0$. If $c_1>c_0$, it reads $(I+\rho T^{-1})Z=T^{-1}Y/c_1$. For
$S\in\{T,T^{-1}\}$, the Neumann identity gives $(I+\rho S)^{-1}=(1-(-\rho)^L)^{-1}\sum_{i<L}(-\rho S)^i$. Since the $L$
powers of $S$ occupy distinct columns in each row, every absolute row sum of $(I+\rho S)^{-1}$ equals
$(1+\rho+\cdots+\rho^{L-1})/|1-(-\rho)^L|$, which is $1/(1-\rho)$ for even $L$ and
$\frac{1-\rho^L}{1+\rho^L}\cdot\frac{1}{1-\rho}\le\frac{1}{1-\rho}$ for odd $L$. Because $c_{\max}(1-\rho)=|c_0-c_1|$,
part~(iv) gives $\|B^{-1}\|_\infty\le 1/|c_0-c_1|\le 4t/7$. Since $B_a$ is a direct sum of such operators on cycles of
even length $A_a$, we have $\|B_a^{-1}\|_\infty\le 4t_a/7\le t_a$.

\emph{(vi) Product.} By submultiplicativity of $\|\cdot\|_\infty$ and part~(v), we obtain
$\|(B_2\cdots B_d)^{-1}\|_\infty=\allowbreak\|B_d^{-1}\cdots B_2^{-1}\|_\infty\le\allowbreak\prod_{a=2}^d(4t_a/7)=\allowbreak(4/7)^{d-1}R/t_1\le\allowbreak R/t_1$.
Telescoping
\[
  \cA_2\cdots\cA_d-B_2\cdots B_d=\sum_{a=2}^d B_2\cdots B_{a-1}(\cA_a-B_a)\cA_{a+1}\cdots\cA_d
\]
and applying the product bound together with parts~(i) and~(ii), $\|F_R x\|_\infty\le\|x\|_\infty$, and
Lemma~\ref{lem:bridge-A} yields~\eqref{eq:lemma-B-total}.
\end{proof}

Notice that the spread array has the same shape $t_1\times\cdots\times t_d$ as the input: the bridge expands no axis.
Appendix~\ref{app:deferred} proves two converse bounds within this two-alias framework. With the positive carrier $e^{+\pi i s/t_a}g_0(s)$ and retained pair
$\nu\in\{0,1\}$, the sample $u=0$ lands at $y=-\frac{1}{2t_a}\notin(0,1)$ and loses uniform $O(2^{-\tw})$ alias decay,
as shown in Proposition~\ref{prop:positive-fails}. Likewise, for any real-valued kernel, $|c_0|=|c_1|$ at the Nyquist
sample $\xi=\tfrac12$, forcing $\|B^{-1}\|_\infty\ge\Omega(L)$ on cycles of length $L=A_a$ by
Proposition~\ref{prop:real-kernel-fails}.

\subsection{The fixed-point bridge algorithm and middle-transform contract}\label{sec:bridge-alg}

This subsection specifies the five-step fixed-point bridge algorithm on dyadic grids and states the interface contract
required of its middle transform.

\paragraph{Parameters.}
Throughout the remainder of this section, $R=t_1\cdots t_d=2^k$ has \emph{balanced shape}: $d\ge 2$, $t_d=r=2^\ell$ with
$\ell\ge 2$, and $t_i\in\{r/2,r\}$ for $1\le i<d$. For the input $x$, we write $X=\max(1,\|x\|_\infty)$. We denote by
$P$ the internal working precision of the bridge, by $q=P+2$ its dyadic grid precision, and by $v=2^{-q}$ the unit of
the $q$-grid $\Cg_q$.

The bridge interacts with the multidimensional transform at its core solely through a black-box accuracy and bit-width
interface. The next definition states this contract: given a bounded spread tensor on the $q$-grid, the middle transform
must return an approximation to $F_{\mathbf t}U$ on the $q$-grid within $E_{\rm mid}v$ in natural frequency order.

\begin{definition}[Middle-transform contract]\label{def:middle-contract}
A \emph{middle transform at precision $q$ with error factor $E_{\rm mid}$ and widening $\delta_{\rm mid}$} is a
procedure that accepts a complete tensor $U$ in natural axis order $(j_1,\dots,j_d)$, with the last axis contiguous,
entries on $\Cg_q$, and $\|U\|_\infty\le 1/8$, and returns an array $\widehat Y$ on $\Cg_q$ in natural frequency order
$(u_1,\dots,u_d)$, with component width at most $q+2+\delta_{\rm mid}$ ($0\le\delta_{\rm mid}\le P$) and
\[
  \|\widehat Y-F_{\mathbf t}U\|_\infty\le E_{\rm mid}\,v,\qquad E_{\rm mid}\le 16R(k+1).
\]
\end{definition}

In our main algorithm, the Walsh middle-transform corollary (Corollary~\ref{cor:walsh-middle} in
Section~\ref{sec:direct-fft}) supplies a middle transform at precision $q$ with $E_{\rm mid}=12\ell$ and
$\delta_{\rm mid}=0$. It applies to every parameter tuple $(q,d,r,K)$ satisfying the hypotheses of the direct tensor FFT
theorem (Theorem~\ref{thm:direct-fft}) at precision~$q$. Lemma~\ref{lem:root-family} verifies those hypotheses along
the root parameter family of Section~\ref{sec:main-alg}, the grid sizes and precisions that the top-level multiplier uses
as functions of~$n$. For comparison, the Bluestein middle transform analyzed in Lemma~\ref{lem:bridge-E}
of Appendix~\ref{app:bluestein} also satisfies Definition~\ref{def:middle-contract}, with
$E_{\rm mid}=\sqrt2\,M(3\ell+1)+12d+16$ and $\delta_{\rm mid}=\lg M$, where $M=R/r$.

Given $x\in\Cg_{p+2}^R$ with $\|x\|_\infty\le 2^p$, the bridge computes $\widehat F_R x$ in five steps:
\begin{enumerate}[label=(\arabic*),leftmargin=*]
  \item \emph{Pre-scaling.} Scan $x$ once to find an integer $H$ with $X\le 2^H<4X$, set
    $S=2^{H+\lceil(\lg(R/t_1))/2\rceil+4}$, and form $x'=x/S$ without rounding on the $q$-grid.

  \item \emph{Spreading.} Reorder the input digits $(j_d,\dots,j_1)$ into natural tensor order $(j_1,\dots,j_d)$ and
    spread the axes $a=d,d-1,\dots,2$ in descending order via~\eqref{eq:spread-axis}, exposing each axis~$a$ as the
    innermost axis by Lemma~\ref{lem:bridge-D}. Along each exposed line of length $t_a$, spreading is a complex cyclic
    convolution with the reflected filter $e\mapsto H^{(a)}_{\delta_a(\mathbf j)}(-e)$ truncated to unfolded offsets
    $|s|\le P$, whose taps are generated once per line.

    We evaluate each line by overlap-add with signed Kronecker packing (cf.~\cite{MS}). First, partition the line into
    blocks of $W_t=\min(2P+1,t_a)$ entries. For each block, pack its signed real and imaginary numerators
    $(u_j)_{j<W_t}$, and those of the $W_t$-tap reflected filter, each into one signed integer $\sum_{j<W_t}u_j B_0^j$
    in base $B_0=2^{4P+16}$ via a one-pass borrow scan across $(4P+16)$-bit slots. Second, compute four exact real integer
    products per block, and extract the exact convolution coefficients sequentially in the balanced range
    $[-B_0/2,B_0/2)$ by subtracting each balanced residue modulo $B_0$ and shifting right by $4P+16$ bits. Third,
    accumulate the overlaps at fractional scale $2q$, wrapping positions modulo $t_a$, and round each completed line
    entry once to $\Cg_q$ by $\Quant_q$.

  \item \emph{Middle transform.} Apply a middle transform at precision $q$ satisfying
    Definition~\ref{def:middle-contract} to the computed spread array $\widehat U$, producing $\widehat Y$.

  \item \emph{Deconvolution.} Apply two-pass cyclic streaming solvers for $B_2^{-1},B_3^{-1},\dots,B_d^{-1}$ in that
    order on the array in natural one-dimensional frequency order, as analyzed in parts~(iv) and~(v) of
    Lemma~\ref{lem:bridge-C} and Section~\ref{sec:bridge-D}.

  \item \emph{Output.} Scale the deconvolved array by the power-of-two factor $S$ without rounding, and round once to
    $\Cg_{p+2}$ by $\Quant_{p+2}$.
\end{enumerate}

All precomputed tables of folded taps, solver parameters, and unit phases are faithful approximations in the sense of
Section~\ref{sec:prelim-model}, certified in part~(iii) of Lemma~\ref{lem:bridge-C}.

\subsection{Component error bounds and the precision contract}\label{sec:bridge-C}\label{sec:bridge-thm}

This subsection bounds the numerical error and bit cost of each finite-precision bridge stage and combines them into the
bridge precision contract (Theorem~\ref{thm:bridge-precision}).

Before assembling the full bridge, we bound each local source of finite-precision error on the $q$-grid ($q=P+2$,
$v=2^{-q}$) for $P\ge 20$, $8\le\tw\le P$, and rational $\lambda_G\in[1,4]$ in $\alpha^2=4\lambda_G\tw\ln 2/\pi$. The
next lemma shows that truncating the unfolded filter to $|e|\le P$ leaves an exponentially small spatial tail. With taps
certified within $2v$, Kronecker-packed spreading of a $q$-grid array of norm at most~$1$ loses at most $4(P+1)v$ per
axis. Rational series evaluate all kernel and phase tables without a transcendental oracle. Moreover, on even cycles with
faithfully stored parameters and $v\le 1/(64t_a)$, the two-pass cyclic solver for $B_a^{-1}$ incurs error at most
$32t_a^2 v\max(1,\|Y\|_\infty)$, independently of the cycle length.

\begin{lemma}[Finite-precision spreading and deconvolution]\label{lem:bridge-C}
Let $P\ge 20$, $q=P+2$, $v=2^{-q}$, and $8\le\tw\le P$. Let $\alpha^2=4\lambda_G\tw\ln 2/\pi$ with a fixed rational
$\lambda_G\in[1,4]$ (so $\alpha$ satisfies~\eqref{eq:alpha}), and let $t=t_a$ ($2\le a\le d$).
\begin{enumerate}[label=(\roman*),leftmargin=*]
  \item \textup{(Filter norms and spatial tail)} For every $\delta\in[0,1)$, $\|H^{(a)}_\delta\|_2\le 1$ and
    $\|H^{(a)}_\delta\|_1\le\sqrt{t}$ over one period $[t]$. Truncating the unfolded kernel to $|e|\le P$ leaves a
    spatial tail $\sum_{|e|>P}|g_a(e+\delta)|<0.10\cdot 2^{-P}=0.40\,v$.
  \item \textup{(Spreading)} With complex taps certified to error at most $2v$, step~(2) of
    Section~\ref{sec:bridge-alg} maps any $q$-grid array $z$ with $\|z\|_\infty\le 1$ to a $q$-grid array within
    $4(P+1)v$ of $\mathcal S_a z$. Each line of length $t$ uses at most $2t/W_t$ blocks ($W_t=\min(2P+1,t)$), with four
    exact products of integers of at most $W_t(4P+16)+1$ bits per block and $O(tP)$ bit operations for packing,
    extraction, and exact accumulation.
  \item \textup{(Tables and phases)} For every integer $b\ge\max(8,P)$, each unfolded kernel sample, alias weight
    $c_0,c_1$, solver parameter $c_{\max}^{-1}$ and $\rho=\min(c_0,c_1)/c_{\max}$, and rational-angle unit phase used by
    the bridge can be computed deterministically with certified error at most $2^{-b}$ in $O((b+k)^3)$ bit operations,
    without any transcendental oracle. In particular, every truncated folded tap can be stored on $\Cg_q$ with certified
    error below $v$ (in $O(P(P+k)^3)$ bit operations per tap), the solver parameters $\hat\mu,\hat\rho$ on $2^{-q}\Z$
    with $\hat\rho\ge 0$ and error below $v$, every unit phase $e$ as a disk-grid entry $\hat e\in\mathbb{D}_q$ with
    $|\hat e-e|<2v$, and a $d$-factor tensor chirp in $\mathbb{D}_q$ with error at most $4dv$.
  \item \textup{(Cyclic solver)} Let $L\ge 2$ be an even cycle length (every bridge cycle length $A_a$, $a\ge 2$, is
    even). Suppose $v\le 1/(64t)$ and the stored parameters satisfy $|\hat\mu-c_{\max}^{-1}|\le v$,
    $|\hat\rho-\rho|\le v$, $\hat\rho\ge 0$, and $|\hat\gamma-(1-\hat\rho^L)^{-1}|\le v$. Then the two-pass cyclic sweep,
    with every product rounded by $\Quant_q$, returns an approximation $\widehat Z$ to $Z=B^{-1}Y$ satisfying
    \[
      \|\widehat Z-Z\|_\infty\le 32t^2\,v\max(1,\|Y\|_\infty),
    \]
    and every intermediate state has modulus strictly below $7t\max(1,\|Y\|_\infty)$.
  \item \textup{(Closure parameter)} Under $v\le 1/(64t)$ and even $L\ge 2$, computing $\hat\rho^L$ from the stored
    $\hat\rho$ by binary powering with truncation at precision $Q=q+\lceil\lg L\rceil+2\lceil\lg(4t)\rceil+10$, followed
    by one guarded division and rounding to $2^{-q}\Z$, produces $\hat\gamma$ with
    $|\hat\gamma-(1-\hat\rho^L)^{-1}|\le v$ using $O(P+\lg L+\lg t)$-bit arithmetic (which is $O(P+k)$ for $L=A_a\le R$).
  \item \textup{(Product norm)} $\|(B_2\cdots B_d)^{-1}\|_\infty\le R/t_1$.
\end{enumerate}
\end{lemma}

\emph{Proof idea.} Part~(i) bounds the filter norms via Parseval's identity and estimates the truncated Gaussian tail by
an exponential integral. Part~(ii) combines this tail with tap approximation error and checks that the Kronecker base
$B_0=2^{4P+16}$ prevents carry wrap during signed overlap-add. Parts~(iv) and~(v) analyze the two-pass cyclic recurrence
for $(I+\hat\rho T)^{-1}$. For even cycle length~$L$, the geometric sum $S_L=\sum_{j<L}\hat\rho^j$ multiplies the
per-step rounding error in the prefix pass and cancels against the closure denominator $1-\hat\rho^L$, yielding a bound
independent of~$L$.

\begin{proof}
\emph{(i) Filter norms and spatial tail.} For $u\in[t]$, unfolding $H^{(a)}_\delta$ and applying~\eqref{eq:poisson}
gives $\sum_{e\in[t]}H^{(a)}_\delta(e)e^{-2\pi i eu/t}=\sum_{\nu\in\Z}G(\frac{u+1/2}{t}-\nu)e^{2\pi i\delta(u/t-\nu)}$,
whose modulus is at most $\sum_{\nu\in\Z}G(\frac{u+1/2}{t}-\nu)=1$. Parseval's identity over $\Z_t$ yields
$\|H^{(a)}_\delta\|_2\le 1$, and Cauchy--Schwarz gives $\|H^{(a)}_\delta\|_1\le\sqrt{t}$. For the spatial tail,
$|g_a(s)|\le e^{-\pi s^2/\alpha^2}/(\pi|s|)$. Since $\alpha\le 4\sqrt{\tw\ln 2/\pi}$ and $\tw\le P$, we have
$\alpha\sqrt{P\ln 2/\pi}\le(4\ln 2/\pi)\sqrt{\tw P}<P$. Thus on each side of $[-P,P]$, the first omitted sample
contributes less than $2^{-P}/(\pi P)$. Substituting $z=\pi s^2/\alpha^2$ and using
$\int_b^\infty e^{-z}z^{-1}\,dz\le e^{-b}/b$ for $b\ge P\ln 2$ bounds the remaining samples by
$\int_P^\infty e^{-\pi s^2/\alpha^2}(\pi s)^{-1}\,ds\le 2^{-P}/(2\pi P\ln 2)$. Summing both sides gives
$\sum_{|e|>P}|g_a(e+\delta)|\le(\frac{2}{\pi P}+\frac{1}{\pi P\ln 2})2^{-P}<0.06\cdot 2^{-P}<0.10\cdot 2^{-P}=0.40\,v$
for $P\ge 20$.

\emph{(ii) Spreading.} Truncating the unfolded kernel to $|s|\le P$ leaves at most $2P+1$ effective taps, folded modulo
$t$ when $t<2P+1$. Each tap is certified to error at most $2v$, while the omitted spatial tail adds at most $0.40\,v$
by part~(i). Hence the effective filter error has $\ell^1$ norm at most $(4P+2.4)v$. Since $\|z\|_\infty\le 1$ and the
single final rounding $\Quant_q$ adds at most $\frac{\sqrt2}{2}v<v$, each output entry differs from $\mathcal S_a z$ by
at most $(4P+3.4)v<4(P+1)v$.

For the packing width, $\|z\|_\infty\le 1$, and part~(i) gives $\|H^{(a)}_\delta\|_\infty\le\|H^{(a)}_\delta\|_2\le 1$,
so every certified tap has modulus below~$2$. At base $B_0=2^{4P+16}$, each numerator satisfies
$|u_j|<2\cdot 2^q<B_0/2$, so $|\sum_{j<W_t}u_j B_0^j|<B_0^{W_t}/2<2^{W_t(4P+16)}$ fits in at most $W_t(4P+16)+1$ signed
bits and is formed by a linear borrow scan in $O(W_t P)$ bit operations. Likewise, every real product coefficient $c_m$
and overlap sum at fractional scale $2q$ has magnitude at most $16(2P+1)2^{2q}<2^{4P+15}=B_0/2$, so
$c_m\in(-B_0/2,B_0/2)$ is the unique balanced residue of $\sum_{k\ge m}c_k B_0^{k-m}$ modulo $B_0$. Reading
$(4P+16)$-bit slots from low to high, subtracting $c_m$, and shifting right by $4P+16$ bits recovers all $c_m$ and
accumulates overlaps at scale $2q$ with no wrap or intermediate rounding. Partitioning the line into
$\lceil t/W_t\rceil\le 2t/W_t$ blocks gives four exact integer products of width at most $W_t(4P+16)+1$ bits per block,
plus $O(tP)$ bit operations for packing, extraction, and accumulation.

\emph{(iii) Tables and phases.} This is Lemma~\ref{lem:phase-tables}, proved in Appendix~\ref{app:deferred}.

\emph{(iv) Cyclic solver.} Assume $c_0>c_1$; the opposite case $c_1>c_0$ reduces to the same forward recurrence via the
reorientation $z_i=Z_{(-i)\bmod L}$ and $y'_i=Y_{(1-i)\bmod L}$, encoding the one-step shift from part~(v) of the proof
of Lemma~\ref{lem:bridge-B}. Let $T$ be the cyclic shift $(TZ)_j=Z_{(j-1)\bmod L}$. Starting from $z=0$, the first pass
computes $b_j=\Quant_q(\hat\mu Y_j)$ and updates $z\leftarrow b_j+\Quant_q(-\hat\rho z)$ for $j=0,\dots,L-1$. It then
closes the cycle via $\widehat Z_{L-1}=\Quant_q(\hat\gamma z)$, and a second pass from $z=\widehat Z_{L-1}$ computes
$\widehat Z_0,\dots,\widehat Z_{L-2}$ by the same recurrence.

By parts~(iii) and~(iv) of Lemma~\ref{lem:bridge-B} and $v\le 1/(64t)$, we have $\rho\le 1-7/(4t)$, which implies $0\le\hat\rho\le 1-1/(2t)$,
$\|(I+\rho T)^{-1}\|_\infty\le t$, and $\|(I+\hat\rho T)^{-1}\|_\infty\le 2t$. Moreover $c_{\max}^{-1}<2.0000001$, so
$|\hat\mu|<2.01$. By the resolvent identity, the exact solution $\tilde Z=\hat\mu(I+\hat\rho T)^{-1}Y$ for the stored
parameters $(\hat\mu,\hat\rho)$ satisfies
$\|\tilde Z-Z\|_\infty\le(2t+2.01\cdot 2t^2)v\|Y\|_\infty\le 6t^2 v\|Y\|_\infty$. Each step of the first pass introduces
at most two nearest roundings of $\frac{\sqrt2}{2}v$ each, adding at most $2v$ in total. Because earlier steps are
damped by powers of $\hat\rho$, the accumulated prefix error at the end of the pass is at most $2v S_L$ with
$S_L=\sum_{j<L}\hat\rho^j$. Since $L$ is even, $(I+\hat\rho T)^{-1}$ closes with denominator
$1-(-\hat\rho)^L=1-\hat\rho^L$, giving the exact cancellation
\[
  \frac{S_L}{1-\hat\rho^L}=\frac{1}{1-\hat\rho}\le 2t.
\]
Relative to $\tilde Z_{L-1}$, the prefix rounding error therefore contributes at most $2v S_L/(1-\hat\rho^L)\le 4tv$ to
$\widehat Z_{L-1}$. Since the prefix states satisfy
$|z|\le S_L(\max_j|b_j|+2v)<4.1t\max(1,\|Y\|_\infty)$, the closure error $|\hat\gamma-(1-\hat\rho^L)^{-1}|\le v$ adds at
most $v|z|<4.1tv\max(1,\|Y\|_\infty)$, and the final $\Quant_q$ adds at most $v$, so
$|\widehat Z_{L-1}-\tilde Z_{L-1}|\le 12tv\max(1,\|Y\|_\infty)$ independently of $L$. The second pass adds at most
$2v/(1-\hat\rho)\le 4tv$, giving
\[
  \|\widehat Z-Z\|_\infty\le(6t^2+16t)v\max(1,\|Y\|_\infty)\le 32t^2 v\max(1,\|Y\|_\infty).
\]
Finally, $\|\tilde Z\|_\infty\le|\hat\mu|\,\|(I+\hat\rho T)^{-1}\|_\infty\|Y\|_\infty<4.02t\|Y\|_\infty$, and every
second-pass state lies within $16tv\max(1,\|Y\|_\infty)\le\tfrac14\max(1,\|Y\|_\infty)$ of $\tilde Z$, so all states have
modulus strictly below $7t\max(1,\|Y\|_\infty)$.

\emph{(v) Closure parameter.} Binary powering of $\hat\rho\in[0,1)$ with truncation toward zero at precision $Q$ incurs
error at most $2L\cdot 2^{-Q}$ and never exceeds $\hat\rho^L$. Since $1-\hat\rho^L\ge 1-\hat\rho\ge 1/(2t)$, both exact
and truncated denominators are at least $1/(2t)$, where the derivative of $x\mapsto(1-x)^{-1}$ is at most $4t^2$. The
powering error therefore perturbs $(1-\hat\rho^L)^{-1}$ by at most $4t^2\cdot 2L\cdot 2^{-Q}\le v/2048$, and one guarded
division followed by nearest rounding to $2^{-q}\Z$ stays within error $v$. Because $\hat\gamma$ is computed from the
stored dyadic $\hat\rho$ without any exact power oracle, the cancellation $S_L/(1-\hat\rho^L)=1/(1-\hat\rho)$ in
part~(iv) holds with no condition relating $L$ to $2^P$.

\emph{(vi) Product norm.} This is part~(vi) of Lemma~\ref{lem:bridge-B}.
\end{proof}

We now combine the stage bounds of Lemma~\ref{lem:bridge-C} with the truncation estimate of Lemma~\ref{lem:bridge-B} to
state the main precision interface of the bridge. For $p\ge\max(8,k)$ and an input $x$ on the $(p+2)$-grid of norm at
most $2^p$, assume the precision schedule~\eqref{eq:schedule} (which makes $P=O(p+\lg R)$), default~$\alpha$, faithful
tables, guarded closure parameters, and any middle transform at precision~$q$ satisfying
Definition~\ref{def:middle-contract}. Then the fixed-point bridge computes $F_R x$ within $\ell^\infty$ error below
$0.18\cdot 2^{-p}\max(1,\|x\|_\infty)$, with the constant made explicit in~\eqref{eq:bridge-total}. Every stored and
intermediate number stays within $O(P)$ bits. Later sections invoke the bridge solely through this schedule~$P$, the grid
precision $q=P+2$, and the component width $w=q+2$.

\begin{theorem}[Bridge precision contract]\label{thm:bridge-precision}
Let $p\ge\max(8,k)$ and let $x\in\Cg_{p+2}^R$ satisfy $\|x\|_\infty\le 2^p$, with $X=\max(1,\|x\|_\infty)$. Set
\begin{equation}\label{eq:schedule}
  \tw=p+k+\lceil\lg d\rceil+20,\qquad
  P=2p+12k+32\lceil\lg(p+k+2)\rceil+256,\qquad
  q=P+2,\qquad
  w=q+2,
\end{equation}
take $\alpha=2\sqrt{\tw\ln 2/\pi}$, certify real table entries to error $v=2^{-q}$ and complex taps and unit phases to
$2v$, and compute closure parameters at precision $Q=q+\lceil\lg L\rceil+2\lceil\lg(4t)\rceil+10$. Then for every middle
transform at precision $q$ satisfying Definition~\ref{def:middle-contract}, the algorithm of
Section~\ref{sec:bridge-alg} returns $\widehat F_R x\in\Cg_{p+2}^R$ satisfying
\begin{equation}\label{eq:bridge-total}
  \|\widehat F_R x-F_R x\|_\infty<\Bigl(\frac{\sqrt2}{8}+2^{-12}+2^{-23}\Bigr)2^{-p}X<0.177021\cdot 2^{-p}X.
\end{equation}
Specifically, before the final rounding the arithmetic error is at most $2^{-p-12}X$ and the omitted-alias error is at
most $2^{-p-23}X$, while the final rounding adds at most $\frac{\sqrt2}{8}\,2^{-p}$. Every stored and intermediate
component has width $O(P)$ bits (Table~\ref{tab:widths}), with $P=O(p+\lg R)$.
\end{theorem}

\emph{Proof idea.} From the linear precision schedule~\eqref{eq:schedule} we first derive a single master margin
inequality~\eqref{eq:master-margin}, which simultaneously discharges all smallness hypotheses of Lemma~\ref{lem:bridge-C}.
Pre-scaling shrinks the input by $\Theta(\sqrt{R/t_1})$ so that every intermediate spread array stays within norm $1/8$,
meeting the input bound of Definition~\ref{def:middle-contract}. Propagating the spreading, middle-transform, and
cyclic-solver errors through $B_2^{-1},\dots,B_d^{-1}$ yields bounds that do not depend on the cycle lengths $A_a$, and
adding the omitted-alias bound and final rounding gives~\eqref{eq:bridge-total}.

\begin{proof}
Write $U=\mathcal S_2\cdots\mathcal S_d x'$ for the exact spread array of $x'=x/S$ and $\widehat U$ for the computed
spread array passed to the middle transform.

\emph{Margins.} Put $z=p+k+2\ge 12$ and $b=\lceil\lg z\rceil$. From~\eqref{eq:schedule}, $P+1\le 44z+233\le 64z$, so
$\lg(P+1)\le b+6$; also $d\le k$. Consequently
$\log_2\bigl(2^{14}d(P+1)R^3 v\bigr)\le 18+4k+b-P\le -p-12$, since $P-(p+4k+b+30)=p+8k+31b+226>0$. This master margin
\begin{equation}\label{eq:master-margin}
  2^{14}d(P+1)R^3 v\le 2^{-p-12}
\end{equation}
implies $16d(P+1)\sqrt{R/t_1}\,v<\frac{1}{16}$, $E_{\rm mid}v\le 16R(k+1)v<\frac{1}{100}$,
$v\le 1/(64\max_a t_a)$, and $64v\sum_{a=2}^d t_a<1$, while $8\le\tw\le P$, verifying all hypotheses of
Lemma~\ref{lem:bridge-C}.

\emph{Pre-scaling.} Let $C=\max(1,\max_j|\mathrm{Re}\,x_j|,\max_j|\mathrm{Im}\,x_j|)$ and $H=1+\lceil\log_2 C\rceil$;
then $X\le 2^H<4X$ and $H\le p+1$. With $S=2^{H+\lceil(\lg(R/t_1))/2\rceil+4}$ we obtain
$16\sqrt{R/t_1}\,X\le S<128\sqrt{R/t_1}\,X$, whence $\|x'\|_\infty\le 1/(16\sqrt{R/t_1})$. Moreover $x'$ has at most
$(p+2)+H+\lceil(\lg(R/t_1))/2\rceil+4\le 2p+\lceil(\lg(R/t_1))/2\rceil+7\le q$ fractional bits, so $x'\in\Cg_q^R$ is
exact.

\emph{Spreading.} By part~(i) of Lemma~\ref{lem:bridge-C}, $\|\mathcal S_a\|_\infty\le\sqrt{t_a}$, so every exact partial
spread has $\ell^\infty$ norm at most $\|x'\|_\infty\prod_{a=2}^d\sqrt{t_a}=\|x'\|_\infty\sqrt{R/t_1}\le 1/16$. By
part~(ii) of Lemma~\ref{lem:bridge-C} and induction over the $d-1$ stages $a=d,\dots,2$, the spreading error
$E_{\rm spread}=\|\widehat U-U\|_\infty$ satisfies
\[
  E_{\rm spread}\le 4(P+1)v\sum_{a=2}^d\prod_{j=2}^{a-1}\sqrt{t_j}\le 16d(P+1)\sqrt{R/t_1}\,v<\frac{1}{16}
\]
by~\eqref{eq:master-margin}. Every computed partial spread therefore has norm at most $\frac{1}{16}+\frac{1}{16}=\frac18\le 1$,
maintaining the input bound of part~(ii) of Lemma~\ref{lem:bridge-C} at each stage and ensuring
$\|\widehat U\|_\infty\le 1/8$ for Definition~\ref{def:middle-contract}.

\emph{Middle transform.} Applying Definition~\ref{def:middle-contract} to $\widehat U$ and using
$\|F_{\mathbf t}\|_\infty\le 1$, $E_{\rm mid}\le 16R(k+1)$, and $k+1\le P+1$, the middle-transform output $\widehat Y$
satisfies $\|\widehat Y-F_{\mathbf t}U\|_\infty\le E_Y:=E_{\rm spread}+E_{\rm mid}v\le 32d(P+1)Rv$ and
$\|\widehat Y\|_\infty<\frac18+\frac{1}{100}<1$.

\emph{Deconvolution.} Put $R_a=\prod_{j=2}^a t_j$ for $2\le a\le d$ and $R_1=1$, so $R_d=R/t_1$. Because
$32v\sum_{a=2}^d t_a<\ln 2$, induction with part~(iv) of Lemma~\ref{lem:bridge-C} and $\|B_a^{-1}\|_\infty\le t_a$ shows
that the input to solver stage~$a$ has norm at most $2R_{a-1}$. Stage~$a$ therefore introduces local error at most
$32t_a^2 v\cdot 2R_{a-1}$, which the subsequent exact inverses $B_d^{-1}\cdots B_{a+1}^{-1}$ amplify by at most
$R_d/R_a$. Since $t_a^2 R_{a-1}(R_d/R_a)=t_a(R/t_1)$, which does not depend on the cycle lengths $A_a$, the total solver
error satisfies
\[
  E_{\rm solver}\le 64\,\frac{R}{t_1}\,v\sum_{a=2}^d t_a\le 64dR^2 v.
\]
Meanwhile, part~(vi) of Lemma~\ref{lem:bridge-C} bounds the amplification of $E_Y$ by
$\|(B_2\cdots B_d)^{-1}\|_\infty\le R/t_1\le R$.
The normalized arithmetic error is thus at most $(R/t_1)E_Y+E_{\rm solver}\le 96d(P+1)R^2 v$, and after exact scaling by
$S<128\sqrt{R/t_1}\,X\le 128RX$,
\[
  E_{\rm arith}\le 2^{14}d(P+1)R^3 v X\le 2^{-p-12}X
\]
by~\eqref{eq:master-margin}.

\emph{Aliases and output.} Applying~\eqref{eq:lemma-B-total} to $x'$ and rescaling by $S$ bounds the omitted-alias error
by $0.120(d-1)(R/t_1)2^{-\tw}X<0.120\cdot 2^{-p-20}X<2^{-p-23}X$, since $(d-1)2^{-\lceil\lg d\rceil}<1$ and $R/t_1<2^k$.
Because $X\ge 1$, nearest rounding $\Quant_{p+2}$ adds at most
$\frac{\sqrt2}{2}2^{-p-2}=\frac{\sqrt2}{8}2^{-p}\le\frac{\sqrt2}{8}2^{-p}X$. Since $1.4142136^2>2$ gives
$\sqrt2<1.4142136$, rational arithmetic yields $\frac{\sqrt2}{8}+2^{-12}+2^{-23}<0.17702096<0.177021$,
proving~\eqref{eq:bridge-total}. All component widths match Table~\ref{tab:widths}; any middle-transform widening
$\delta_{\rm mid}$ is stripped only after verifying $\|\widehat Y\|_\infty<1$, which alters no value.
\end{proof}

\begin{table}[ht]
\centering
\small
\caption{Bit widths of all stored and intermediate quantities inside the fixed-point bridge. Widths follow the precision
schedule of Theorem~\ref{thm:bridge-precision} at grid precision $q=P+2$; on the direct tensor FFT route of
Section~\ref{sec:direct-fft}, the middle-transform widening is $\delta_{\rm mid}=0$.}\label{tab:widths}
\begin{tabular}{@{}>{\raggedright}p{0.38\textwidth}>{\raggedright\arraybackslash}p{0.58\textwidth}@{}}
\toprule
Object & Width or scale\\
\midrule
normalized input $x'$, completed spread arrays & $q+2=P+4$ bits per real/imaginary component\\
table evaluation (Lemma~\ref{lem:phase-tables}) & $S=2b+8\lceil\lg b\rceil+64$ fractional bits ($b\le q+O(\lg P)$) plus $O(b+k)$ integer bits; temporaries $\le 3S+O(k)$ fractional bits\\
exact overlap-add accumulators & fractional scale $2q$, plus $O(\lg P)$ integer bits\\
packed spread operands & $O(P)$ slots of $4P+16$ bits\\
middle-transform output $\widehat Y$ & $q+2+\delta_{\rm mid}$ bits ($0\le\delta_{\rm mid}\le P$)\\
deconvolution states and prefixes & $q+k+6=P+k+8$ bits\\
closure powering (Lemma~\ref{lem:bridge-C}(v)) & $Q=q+\lceil\lg L\rceil+2\lceil\lg(4t)\rceil+10$ fractional, $O(\lg t)$ integer bits\\
final exact scaling by $S$ and output & add $\log_2 S$ to exponent; one rounding $\Quant_{p+2}$ to $\Cg_{p+2}$\\
\bottomrule
\end{tabular}
\end{table}

\subsection{Axis exposures and streaming layout}\label{sec:bridge-D}

This subsection shows how to expose each spread axis on a multitape Turing machine and stream the cyclic deconvolutions
without full array transpositions.

Spreading operates on the lines of one axis~$a$ at a time, which lie contiguously on the tape only when that axis is
innermost. Bubbling each axis inward through $O(d)$ adjacent swaps would cost $O(Vd^2\ell^\tau)$ steps per sweep over
the $d$ axes, more than the cost analysis of Section~\ref{sec:main} can absorb at our dimension exponent~$\eps$.
Instead, we exchange axis~$a$ directly with the innermost axis in a single call to the arbitrary-width interchange
(Lemma~\ref{lem:transfer12}). The Walsh positional
permutations of Section~\ref{sec:walsh} do not apply here because bridge line operations carry $O(P)$-bit scalar entries
rather than length-$r$ polynomial payloads.

To run a line operation along axis~$a$, we swap axis~$a$ directly with axis~$d$ and swap it back afterwards. The next
lemma bounds the tape cost of each exposure and restoration by $O(V(1+\ell^\tau))$ steps even when the two axis
bit-widths differ by one bit, so sweeping across all $d$ axes costs $O(Vd(1+\ell^\tau))$ steps rather than
$O(Vd^2\ell^\tau)$.

\begin{lemma}[Axis exposures via chunk interchanges]\label{lem:bridge-D}
Let $\tau\in(0,1)$ be any rational interchange exponent satisfying $s_b/W_b<m^\tau$ in Lemma~\ref{lem:transfer12} (for
instance $\tau=1-2^{-40}$ by Proposition~\ref{prop:deficits}, or $\tau=1-2^{-50}$ as in~\cite[Lemma~12]{OAI26}). Let a
complete tensor of logical bit volume $V$ have axis bit-widths in $\{\ell-1,\ell\}$ ($\ell\ge 2$) and a nonempty record
payload. Given the selected axis and the lengths of the collapsed prefix, gap, and suffix fields as a fixed number of
binary integers (their computation from the axis list is charged to the caller), any axis can be exposed as the
innermost axis, and the original axis order restored afterwards, in $O(V(1+\ell^\tau))$ steps. A sweep of line
operations over all axes uses at most $2(d-1)$ such exchanges, and reversing the axis order uses at most
$\lfloor d/2\rfloor$ more, so a constant number of sweeps costs $O(Vd(1+\ell^\tau))$ steps.
\end{lemma}

\emph{Proof idea.} We collapse the axes before the chosen axis~$a$ into a prefix field, the axes strictly between~$a$ and
the innermost axis~$d$ into a gap field, and the record payload into a suffix field, as Lemma~\ref{lem:transfer12}
requires. When the two axes have equal bit-width, one interchange suffices; when their widths differ by one bit, we split
off the extra bit and combine one equal-width interchange with one linear-time single-bit move.

\begin{proof}
Collapse all axes preceding the chosen axis~$a$ into the prefix field $A$, all axes strictly between~$a$ and the
innermost axis~$d$ into the gap field $G$, and the scalar payload into the suffix $B$ of Lemma~\ref{lem:transfer12}. If
axes $a$ and $d$ have equal bit-width $\ell$ or $\ell-1$, one call to Lemma~\ref{lem:transfer12} exchanges them in
$O(V\ell^\tau)$ steps, independently of the size of $G$.

Now suppose their widths differ by one bit. If the earlier field is longer, write it as $(x_0,X)$ with $|X|=|Y|$, where
$Y$ is the later field, and execute $(x_0,X,G,Y)\mapsto(x_0,Y,G,X)\mapsto(Y,G,x_0,X)$. This uses one equal-width
interchange of cost $O(V\ell^\tau)$ followed by one single-bit move of cost $O(V)$ from item~(1) of
Fact~\ref{fact:oai-streams} and item~(2) of Lemma~\ref{lem:transfer-streams}, totaling $O(V(1+\ell^\tau))$. If the later
field is longer, execute $(X,G,y_0,Y)\mapsto(y_0,X,G,Y)\mapsto(y_0,Y,G,X)$. Running the inverse sequence after the line
operation restores the layout. Axis~$d$ requires no exchange, so a full sweep uses at most $2(d-1)$ exchanges, and
reversing the $d$ axes factors into $\lfloor d/2\rfloor$ disjoint exchanges of the same form.
\end{proof}

Two streaming observations complete the tape implementation of Section~\ref{sec:bridge-alg}:
\begin{itemize}[leftmargin=*]
  \item \emph{Deconvolution without transposition.} In natural one-dimensional frequency order, write any frequency
    $f\in[R]$ as $f=j\,(R/A_a)+s$ with $0\le j<A_a$ and $0\le s<R/A_a$. For each fixed $s$, the entries
    $j=0,1,\dots,A_a-1$ form a single cycle of $B_a$ along which $u_a(f)$ is constant, while the tape stores, for
    $j=0,1,\dots,A_a-1$ in order, the contiguous row of all $s\in[R/A_a]$.

    On a sequential state tape, we maintain one $O(P)$-bit solver state record for each $s\in[R/A_a]$, holding the
    current state $z$ and the column parameters $\hat\mu,\hat\rho,\hat\gamma$. Because these parameters depend only on
    $a$ and $u_a$, they are replayed rather than recomputed. We scan each data row $j$ in lockstep with the state tape
    and rewind only the state tape after each row. This costs $O(V)$ tape movement per pass and per axis~$a$,
    independently of $A_a$, while the $O(1)$ multiplications per entry are charged in Section~\ref{sec:main-tables}.
    Columns with $c_0>c_1$ and $c_1>c_0$ are processed in two masked passes. For $c_1>c_0$, linear-time row-block
    reversals implement the index reflections $j\mapsto(-j)\bmod A_a$ and $j\mapsto(1-j)\bmod A_a$ from the proof of
    part~(iv) of Lemma~\ref{lem:bridge-C} and undo them afterwards. This lockstep scan relies on
    Definition~\ref{def:middle-contract} returning $\widehat Y$ in natural frequency order.

  \item \emph{Amortized tap and table generation.} Each exposed axis~$a$ has $R/t_a\le 2R/r$ lines, and generating the
    $O(P)$ folded taps of a line takes $\mathrm{poly}(P,k)$ bit operations by part~(iii) of Lemma~\ref{lem:bridge-C}. Even
    regenerating taps for every line therefore costs $O(dR(P+k)^{O(1)}/r)$ steps, which is $o(RP)$ whenever $r$
    exceeds every fixed polynomial in~$P$ for large~$P$, as holds along the root family of Lemma~\ref{lem:root-family}.
    The same bound applies to the solver and phase tables.
\end{itemize}
